% Part: first-order-logic
% Chapter: completeness
% Section: downward-ls

\documentclass[../../../include/open-logic-section]{subfiles}

\begin{document}

\olfileid{fol}{com}{dls}
\olsection{లొవెన్‌హైమ్--స్కోలెమ్ సిద్ధాంతం}

ఒక సిద్ధాంతానికి అనంత నమూనా ఉంటే, దానికి గరిష్ఠంగా !!{denumerable}
అయిన నమూనా కూడా ఉంటుందని లొవెన్‌హైమ్--స్కోలెమ్ సిద్ధాంతం చెబుతుంది.
ఈ వాస్తవపు తక్షణ పరిణామం ఏమిటంటే, ఒక !!a{structure} పరిమాణం
!!{nonenumerable} అని ప్రథమక్రమ తర్కం వ్యక్తం చేయలేదు: అన్ని
!!{nonenumerable} !!{structure}sలో సంతృప్తిపరచబడే ఏ !!{sentence} లేదా
!!{sentence}s సమితైనా, ఏదో ఒక !!{enumerable} నిర్మాణంలో కూడా
సంతృప్తిపరచబడుతుంది.

\begin{thm}
\ollabel{thm:downward-ls} $\Gamma$ అవైరుధ్యమైతే దానికి ఒక
!!a{enumerable} నమూనా ఉంటుంది; అంటే, పరిమిత లేదా !!{denumerable}
ప్రదేశం గల !!a{structure}లో అది సంతృప్తిపరచదగినది.
\end{thm}

\begin{proof}
$\Gamma$ అవైరుధ్యమైతే, సంపూర్ణతా సిద్ధాంతపు నిరూపణ ఇచ్చే
!!{structure}~$\Struct M$ యొక్క ప్రదేశం~$\Domain{M}$, భాష~$\Lang L$
పదాల సమితికంటే పెద్దది కాదు. కాబట్టి $\Struct M$ గరిష్ఠంగా
!!{denumerable}.
\end{proof}

\begin{thm}
\ollabel{noidentity-ls} తాదాత్మ్యం లేని ప్రథమక్రమ తర్క భాషలోని
!!{sentence}s అవైరుధ్య సమితి~$\Gamma$కు ఒక !!a{denumerable} నమూనా
ఉంటుంది; అంటే, అనంతమైన, !!{enumerable} ప్రదేశం గల !!a{structure}లో
అది సంతృప్తిపరచదగినది.
\end{thm}

\begin{proof}
$\Gamma$ అవైరుధ్యమై, తాదాత్మ్యం సంభవించే వాక్యాలేవీ దానిలో లేకపోతే,
సంపూర్ణతా సిద్ధాంతపు నిరూపణ ఇచ్చే !!{structure}~$\Struct M$ యొక్క
ప్రదేశం~$\Domain{M}$, భాష~$\Lang L'$ పదాల సమితికి తాదాత్మ్యంగా ఉంటుంది.
కాబట్టి $\Trm[L']$లాగే $\Struct{M}$ కూడా !!{denumerable}.
\end{proof}

\begin{ex}[స్కోలెమ్ విరోధాభాసం]
జెర్మెలో--ఫ్రెంకెల్ సమితి సిద్ధాంతం~$\Log{ZFC}$, సమితుల పరిమాణాల
గురించిన వాస్తవాలతో సహా, దాదాపు అన్ని గణిత ప్రకటనలనూ వ్యక్తం చేయగల
చాలా శక్తివంతమైన చట్రం. ఉదాహరణకు, వాస్తవ సంఖ్యల సమితి~$\Real$
!!{nonenumerable} అని~$\Log{ZFC}$ నిరూపించగలదు; ఏ సమితి ఘాత సమితైనా
ఆ సమితికంటే పెద్దదనే కాంటర్ సిద్ధాంతాన్ని కూడా నిరూపించగలదు; ఇంకా
ఇలాంటివే మరెన్నో. $\Log{ZFC}$ అవైరుధ్యమైతే దాని నమూనాలన్నీ అనంతమైనవి;
అంతేకాక, సహజ సంఖ్యల ఘాత సమితి ఉందనే~$\Log{ZFC}$ సిద్ధాంతాన్ని సత్యం
చేసే మూలకం వంటి, !!{nonenumerable} అని సిద్ధాంతం చెప్పే !!{element}sను
అవన్నీ కలిగి ఉంటాయి. లొవెన్‌హైమ్--స్కోలెమ్ సిద్ధాంతం ప్రకారం,
$\Log{ZFC}$కు !!{enumerable} నమూనాలు కూడా ఉన్నాయి---``!!{nonenumerable}''
సమితులను కలిగి ఉండి, తామే !!{enumerable} అయిన నమూనాలు.
\end{ex}

\end{document}
